\clearpage
\fancyhead[L]{\small CARLINI. TAVOLE ASTRONOMICHE.}
\hypertarget{native-548-lower}{}\label{native-548-lower}
\begin{center}\small Göttingische gelehrte Anzeigen. No. 193. pp. 1921..1928. December 3, 1812.\\[3pt]\rule{.72\linewidth}{.3pt}\end{center}

\emph{Esposizione di un nuovo metodo di costruire le tavole astronomiche applicuto alle tavole del sole di} \textsc{Francesco Carlini}. 1810. \emph{Dalla reale stamperia.} LXI and 92 pages in small quarto. Milan.

Thanks to the work of \textsc{Laplace}, \textsc{Delambre} and \textsc{von Zach}, our latest solar tables have reached a degree of accuracy that leaves little to be desired in that respect. Greater \emph{ease of use} was still desirable in calculating positions of the Sun, a task that astronomers have to do so often and that takes up so much of their time. The chief aim of \textsc{von Zach}'s tables was to fit them into a small space, with less attention to ease of use. \textsc{Delambre} took somewhat more care to make his tables easy to use. He supplied the mean motions for every day of the year, a very detailed table for the equation of the center, and double-entry tables for the perturbations due to the other planets. Even so, these tables leave much to be desired. Without any major changes, they would have been easier to use if their author had chosen to recast the perturbations due to the Moon beforehand, to leave several small tables in their proper places instead of moving them out of a very ill-timed economy, which forces one to leaf back and forth every time, and to turn the perturbations of the radius vector into perturbations of its logarithm. On top of this, the perturbations of the radius vector due to Jupiter are calculated incorrectly and from an incomplete formula, so one must first add a correction table to one's own copy.

\SourceLeaf{549}We are pleased to announce to astronomers that this work by \textsc{Carlini} offers new solar tables whose ease of use surpasses everything we have had so far, and indeed everything we have even wished for. One might say that with consummate skill their author has done everything possible to make them easier to use. He has anticipated almost every wish. Calculating a position of the Sun has become a simple game in which one hardly has to write a single digit beyond what is absolutely necessary. We cannot help giving a somewhat fuller account of a work for which astronomers who carry out these calculations owe its author the greatest thanks.

The main idea that \textsc{Carlini} has used so successfully is to let the angular arguments on which a table depends advance in \emph{days}, in place of the degrees or thousandths of a revolution used previously. Thus the equation of the center of the Sun for a given eccentricity (\textsc{Carlini} chose the value for 1810) can be given in a table whose argument is the number of days elapsed since the last passage through perigee. Since we will have to refer to the time of such a passage several more times, we will denote it by $T$ for brevity. The argument is very easy to obtain if one knows how many days had elapsed since the last $T$ or still had to elapse before the next $T$ at the start of each year (\textsc{Carlini} always takes this to be noon on the preceding December 31 at the Milan meridian, including in leap years). For each year this time interval, which may be taken as positive in the first case and negative in the second, can be entered in a separate table as a kind of epact. It is represented by a small number because $T$ happens always to fall close to the start of the year. To make the number always positive, \textsc{Carlini} increased it by two units and entered it in that form in Table II under the heading A. So a date $g$ days after the start of a year falls $g+A-2$ days after the next $T$. The corresponding equation of the center, already combined with the mean motion over those $g+A-2$ days, is found next to the argument $g+A$ in Table III under the heading Prima equazione della longitudine del Sole. To obtain the elliptic longitude of the Sun at once, one need only add the longitude of perigee at time $T$, with that longitude referred to the start of the given year. This longitude is also included in Table II as the Termine constante della longitudine del Sole. The foregoing would fully explain the two headings just mentioned in Tables III and II if the Sun's motion had no other inequalities and the eccentricity were invariable. Since this is not so, the two entries undergo some further modifications, which we will now explain. The secular change of the eccentricity requires a correction to the equation of the center that can be expressed by the formula $(t+f-1810)c$, where $t$ is the year number, $f$ is the fraction of the year, and $c$ is a quantity depending on the mean anomaly that lies between the limits $+0''173$ and $-0''173$. Thus, on setting $c = c^1 - 0''173$, $c^1$ is always positive. The entire correction is then \[ = (t-1810)c^1 - (t-1810)0''173 + fc \] The part $-(t-1810)0''173$ is already combined with the longitude of perigee and included in the Termine constante of Table II. The part $fc$ is combined with the Prima equazione of Table III. Thus the only remaining term to add is $(t-1810)c^1$. For this purpose the always positive $c^1$ is found in Table IV under the heading Variazione annua della prima equazione, with the argument $g+A$ of the third table. The Prima equazione in Table III also includes the part of the nutation that depends on the longitude of the Sun. The table itself goes \SourceLeaf{550}in steps of one tenth of a day up to 370, so interpolation is very easy, since the first difference does not exceed $367''$ and the second vanishes entirely.

We have yet to finish explaining the two headings of Table II. \textsc{Carlini} has a very ingenious way of calculating aberration while taking its variable part into account. In the existing tables the calculated longitude of the Sun includes the constant part of this aberration, $-20''25$. Thus \emph{both} the epoch of the mean longitude and the longitude of perigee are in effect already reduced by $20''25$. To find the apparent longitude of the Sun, one then still has to add the variable part of the aberration from a separate table. As \textsc{Carlini} points out, one could have saved that last operation by reducing the longitude of perigee by only half as much, namely $10''12$, while retaining the longitude of the mean longitude as reduced by $20''25$. For this reason \textsc{Carlini} had to increase the numbers in the column of Table II headed Termine constante by $10''1$ and decrease the numbers $A$ in the same table by 0,00286 day. One then obtains the Sun's longitude with aberration directly. If the true longitude of the Sun is wanted, one need only add $10''1$ to the Termine constante and 0,00286 day to $A$. Finally, the Termine constante also includes the quantity $-56''6$ to offset the sum of all the constants by which \textsc{Carlini} makes positive the perturbation equations we are about to examine more closely. That sum is $+56''6$.

The tables for the small inequalities in the Sun's longitude are based on principles like those of the table for the equation of the center. A suitable transformation reduces the inequalities due to the Moon to equations with single arguments. Of these, \textsc{Carlini} has retained the two largest, whose maxima are $7''5$ and $0''5$. In these tables as well as in the table for nutation, the argument is expressed in days, which actually represent the interval between the day for which one is calculating and the day on which the argument of each table was 0. Table II also gives the value of each argument expressed in this way at the beginning of every year from 1750 to 1900.

Finally, the tables of planetary perturbations are exceptionally easy to use. \textsc{Carlini} calculated them with two arguments, which, as is well known, represent the mean longitude of the Earth and the planet's mean heliocentric elongation from the Earth. Both can again be expressed in days. But \textsc{Carlini} has arranged the table so that, for the first argument, one need only enter the number of days elapsed since the start of the year and, for the second, the value that it had at the start of the year. This value is also given in Table II for the perturbations due to each of the planets Venus, Jupiter, Mars and Saturn. The second argument thus stays the same throughout the year. The best course is to take from the tables the equations for the immediately preceding and following days, which are given directly, in place of those for the day for which one is calculating (so the tables become single-entry tables), and then interpolate between the sums of all four perturbations. Let us make a few minor points about these excellently arranged perturbation tables. \textsc{Carlini} calculated them on the assumption that the mean longitude of the Earth at the start of the year was $= 100^\circ$. What this actually means is that, \emph{if one wants to be completely strict,} one must enter them 1) with the number of days elapsed since the moment at which the mean longitude of the Sun was $= 100^\circ$, in place of the days elapsed since the start of the year, and 2) with the value that the mean elongation of the perturbing planet had at that very moment, in place of its value at the start of the year. The difference is admittedly never appreciable, \SourceLeaf{551}but one might still wish to be able to choose the exact arguments whenever one pleased. For the first argument this could have been done with a separate correction table, like the one we have added to our copy for our own use. The second argument could have been entered in Table II just as easily for the exact moment as for the start of the year. In 1804, for example, the correction to the first argument is 1.09, and the exact values of $E$, $F$, $G$, $H$ are \[ \begin{array}{r|l@{\;}l} E = 36.66 & \text{in the table} & E = 36.55\\ F = 28.45 & & F = 28.34\\ G = 13.13 & & G = 13.09\\ H = 10.1 & & H = 10.1 \end{array} \]

Second, we would have liked the table itself to extend another 10 days, to 370 days instead of 360. When one wants to calculate a position of the Sun for the last days of December, this would save the (admittedly small) trouble of writing out the arguments $E$, $F$, $G$, $H$ for the following year as well.

To give our readers an idea of how easy these tables are to use, we have described in some detail how the ones for calculating longitude are arranged. Our limited space does not allow us to go into the same detail for the other tables. We will therefore remark only that the tables for directly calculating the logarithm of the \emph{distance} of the Sun from the Earth, and the tables for calculating the Sun's \emph{latitude}, are arranged just as deftly as those for longitude. A few tables are also appended for reduction from the ecliptic to the equator, for the mean right ascension of the Sun expressed in time, and for reducing the positions of the Sun from mean noon to true noon.

The introduction that precedes the tables explains how they are arranged with clarity and elegance. It is especially valuable for the many ingenious devices it presents to make calculating an ephemeris of the Sun's positions for an entire year as easy as possible. We do not have room to go into their details here, but we consider it our duty to urge all astronomers who have many solar positions to calculate to study and use these advantages.
